\begin{figure}[H]
\centering
\begin{tikzpicture}[
  >=Latex,
  every node/.style={font=\small},
  state/.style={
    draw, rounded corners=2pt, align=center, inner sep=5pt,
    text width=2.20cm, minimum height=0.82cm
  },
  eqbox/.style={
    draw, rounded corners=2pt, align=center, inner sep=6pt,
    text width=10.8cm, minimum height=1.12cm
  },
  flow/.style={->,thick},
  reverse/.style={->,thick,dashed}
]
\node[font=\normalsize\bfseries] at (6.1,3.25)
  {得分提供反转扩散所需的修正};

% Forward row
\node[state] (f0) at (0,1.75) {$p_0$\\数据};
\node[state] (ft) at (6.1,1.75) {$p_t$\\中间状态};
\node[state] (fT) at (12.2,1.75) {$p_T$\\简单噪声};

\draw[flow] (f0) -- node[above] {\scriptsize 前向加噪} (ft);
\draw[flow] (ft) -- (fT);

\node[eqbox] (forward) at (6.1,0.15) {
\textbf{前向 SDE}\\[3pt]
$\displaystyle
dX_t=f(X_t,t)\,dt+g(t)\,dW_t
$
\qquad
\textbf{沿其边缘分布学习：}
$\displaystyle
s_\theta(x,t)\approx\nabla_x\log p_t(x)
$
};

% Reverse row
\node[state] (r0) at (0,-1.65) {$p_0$\\数据};
\node[state] (rt) at (6.1,-1.65) {$p_t$\\中间状态};
\node[state] (rT) at (12.2,-1.65) {$p_T$\\简单噪声};

\draw[reverse] (rT) -- node[above] {\scriptsize 反向时间生成} (rt);
\draw[reverse] (rt) -- (r0);

\node[eqbox] (reversebox) at (6.1,-3.45) {
\textbf{反向时间 SDE：从 $T$ 向 $0$ 模拟}\\[3pt]
$\displaystyle
dX_t=
\left[
f(X_t,t)-g(t)^2s_\theta(X_t,t)
\right]dt
+
g(t)\,d\bar W_t
$
};

\draw[flow] (forward) -- (reversebox);

\end{tikzpicture}
\caption{\textbf{得分提供反转扩散所需的修正。}
前向 SDE 定义带时间索引的边缘分布 $p_t$，而以时间为条件的模型学习 $s_\theta(x,t)\approx\nabla_x\log p_t(x)$。对于这里采用的标量、与状态无关的扩散系数，反向漂移与前向漂移相差 $-g(t)^2\nabla_x\log p_t(x)$。用 $s_\theta(x,t)$ 替代未知得分，就得到实际使用的反向时间 SDE，并从 $T$ 向 $0$ 进行模拟。}
\label{fig:score-forward-reverse-sde}
\end{figure}